\indent 低空飞行所处的大气环境具有局地变化快、垂直结构复杂和观测尺度不一致的特点。围绕湍流强度评估、三维场重构与航路优化三个相互衔接的问题，本文以多源观测为基础，将物理诊断、统计标定和图搜索相结合，构建从观测参数到飞行风险的计算框架。

\textbf{问题一：}针对风廓线雷达与微波辐射计在观测要素上的互补关系，首先完成时间匹配、高度对齐与质量控制，计算位温、垂直风切变和理查逊数。通过融合热力稳定性与雷达谱宽代理量，建立综合评估模型 a；随后仅使用雷达变量及其历史变化构建模型 b。以模型 a 的输出为参照，采用约束回归与随机森林两条路线标定模型 b，并按时间块比较垂直廓线和误差分布。测试集上的 RMSE、MAE 和决定系数分别为\slot{数值}、\slot{数值}和\slot{数值}，相对未标定模型的误差变化为\slot{数值}。该比较反映模型 b 对参照指标的逼近程度，独立湍流观测用于检验指标的物理有效性。

在参数计算中，位温采用 $\theta=T(p_0/p)^{R_d/c_p}$，以消除气压差异对热力状态比较的影响；风切变采用 $S=[(\partial u/\partial z)^2+(\partial v/\partial z)^2]^{1/2}$，由此构建 $Ri=(g/\theta)(\partial\theta/\partial z)/S^2$。对谱宽完成可获得的非湍流展宽校正后，定义代理量 $q_{SW}$，并建立综合指标 $I_a=\alpha_1f(Ri)+\alpha_2\widetilde q_{SW}+\alpha_3\widetilde S$。其中，稳定度映射与归一化参数在标定阶段固定。模型 b 以雷达当前量和历史变化为特征，通过正则化目标控制拟合误差和系数幅度，再用非线性候选检验变量交互的作用。分高度残差用于判断优化主要发生在哪些层次，参数扰动用于分析参照指标的不确定性。

\textbf{问题二：}围绕多种设备探测方式和空间覆盖不同的问题，将地面自动站、风廓线雷达及多普勒天气雷达转换到统一的时间和坐标系统。利用各向异性距离描述水平与垂直相关尺度，以设备误差、空间距离和时间衰减确定融合权重；进一步建立包含观测项、背景项和平滑项的变分模型，求解三维湍流强度场。输出覆盖全部地面站所在矩形区域，水平网格间距为 100 m、垂直网格间距为 50 m，高度低于 2 km。通过留站验证、设备消融及分高度统计，计算重构误差与不确定性，得到\slot{主要空间特征}及\slot{可靠区域与误差范围}。

数据融合以 $J=J_{\rm obs}+\lambda_bJ_b+\lambda_sJ_s$ 为优化目标，分别描述观测拟合、背景约束和空间平滑。观测算子将网格场映射到点位、垂直门或雷达采样体积，使不同设备在各自代表尺度上参与约束。对于线性算子和二次目标，采用共轭梯度法求解相应正定系统；引入指标范围约束时，使用匹配的有界优化方法。结合时间更新与半拉格朗日平流，描述历史场的连续变化，并通过分块重采样生成误差区间。空间展示同时给出风险、有效观测覆盖和不确定性，以区分直接观测支持与背景约束主导的区域。

\textbf{问题三：}针对有数值预报和无数值预报两类场景，分别建立模型 d 与模型 e。模型 d 利用 02:00 至 05:00 的重构场标定数值天气预报诊断量，并生成 05:00 至 08:00 的预报；模型 e 利用截至 05:00 的历史观测建立非线性外推关系，生成 05:00 至 06:00 的预测。以分时效回报检验比较误差、事件识别率和区间覆盖率，再将预报场转换为时空图上的边代价。在给定起终点和允许飞行区域内，以累计湍流暴露量为优化目标搜索航路，得到\slot{航路编号及风险值}。与同一预报场下的基准航路相比，累计暴露量变化为\slot{数值}，路径长度为\slot{数值}。

在航路层，将位置和到达时间共同表示为图节点，边代价定义为 $c_{ij}=\ell_{ij}[I(n_i)+I(n_j)]/2$。该式把局地风险转化为沿线累计暴露量，路径最优性在给定离散图及允许区域内讨论。A* 启发函数采用单位距离风险下界与剩余距离的乘积；下界为零时使用 Dijkstra 搜索。最终路径由独立程序重新计算边段可行性和风险积分，并通过网格加密与预报情景扰动，比较路径选择对离散尺度和场误差的敏感程度。

上述建模过程将单站诊断、区域重构和路径选择联系起来，各阶段的输出分别对应后续模型的输入与验证参照。最终交付包括垂直廓线、三维网格场、分时效预报、航路节点序列及可运行程序，具体结果以对应计算记录为准。

\noindent\textbf{关键词：}湍流强度量化\quad 多源数据融合\quad 三维湍流场重构\quad 低空航路规划\quad 变分同化模型
